%!TEX root = ../dokumentation.tex

\pagestyle{empty}

\renewcommand{\abstractname}{Abstract} % Zusammenfassung in Abstract umbenennen

\begin{abstract}

    \todo[inline]{Sperrvermerk und Erklärung auch ins Dokument? To-Do: Nochmal nachlesen oder klären}

    Die Arbeit entwickelt ein KI-basiertes Assistenzsystem zur quantitativen Portfoliooptimierung, das prozentuale Handlungsempfehlungen für Einzelwerte mit transparenten, natürlichsprachlichen Begründungen verknüpft und zu einer risikoadaptiven Portfoliosicht einschließlich Rebalancing-Implikationen aggregiert.
    Daraus leiten sich zwei Forschungsfragen zur konstruktiven Realisierung eines solchen Systems sowie zu seiner Fähigkeit ab, eine übergeordnete Risiko-Rendite-Bewertung auf Portfolioebene bereitzustellen.

    Methodisch folgt die Arbeit einem Design-Science-Research-Ansatz und realisiert einen modularen Prototypen mit klaren HTTP/JSON-Schnittstellen: Python/Flask als Backend, eine austauschbare LLM-Schicht (lokal via Ollama, optional per API zu externen Anbietern), \texttt{yfinance} als zentrale Kursdatenquelle, LangChain zur Orchestrierung sowie ein leichtgewichtiges Web-Frontend (HTML/Bootstrap/JS).
    Die klare Projektstruktur und Entkopplung der Komponenten unterstützen Erweiterbarkeit, Austauschbarkeit und Reproduzierbarkeit.

    Die Evaluation ist als ressourcensensitives Backtesting über repräsentative Marktphasen angelegt und kombiniert finanzielle Performancekennzahlen (u.\,a. Sharpe Ratio, Volatilität, maximaler Drawdown, VaR/ES) mit statistischen Absicherungen (t-/Wilcoxon-Tests, Bootstrap) sowie XAI- und Nutzerstudien-Bausteinen.
    Damit wird ein Evaluationsraster bereitgestellt, das Transparenz- und Leistungsaspekte integriert und für zukünftige Auswertungen operationalisiert.

    Die wesentlichen Limitationen betreffen die probabilistische Natur großer Sprachmodelle und deren Prompt-Sensitivität, potenzielle Datenqualitäts- und Abdeckungsgrenzen sowie Rechenkosten, die die Tiefe der Tests begrenzen und die externe Validität einschränken.
    Praktisch adressiert der Prototyp Transparenz-, Dokumentations- und Governance-Anforderungen und schafft eine belastbare Grundlage für vertiefte Evaluationen, funktionsseitige Erweiterungen und anwendungsnahe Pilotierungen.
\end{abstract}


